

\newglossaryentry{FRBR}{
		name={FRBR},
		description={Functional Requirements for Bibliographic Records. Modèle conceptuel entité-association développé par l'IFLA.}
	}

\newglossaryentry{fiaf}{
	name={FIAF},
	description={Fédération Internationale des archives du film. Organisation regroupant les cinémathèques et édité le manuel de catalogage audiovisuel fondé sur le modèle FRBR.}
}


\newglossaryentry{rush}{
	name={rush},
	description={Ensemble des épreuves brutes de tournage, non montées et non sélectionnées, constituant le matériau source d'un film.}
}

\newglossaryentry{Oeuvre}{
	name={Oeuvre},
	description={}
}

\newglossaryentry{Expression}{
	name={Expression},
	description={}
}

\newglossaryentry{Manifestation}{
	name={Manifestation},
	description={}
}

\newglossaryentry{Item}{
	name={Item},
	description={}
}

\newglossaryentry{Archivesaudiovisuelles}{
	name={Archives audiovisuelles},
	description={}
}


\newacronym{WVMI}{WVMI}{Work, Variant, Manifestation, Item.}
\newacronym{IFLA}{IFLA}{International Federation of Library Associations and Institutions}
\newacronym{FRAD}{FRAD}{Fonctionnalités requises des données d'autorité}
\newacronym{FRSAD}{FRSAD}{Fonctionnalités requises des données d'autorité matière}
\newacronym{IFLA-LRM}{IFLA-LRM}{long}
\newacronym{FRBRoo}{FRBRoo}{long}
\newacronym{CIDOC-CRM}{CIDOC-CRM}{object-oriented Conceptual Reference Model}